\section{Introduction}\label{sec:introduction}

Handwriting is a widely studied behavioral biometric, relevant to forensic document examination, signature authentication, and manuscript attribution. As a learned motor skill rather than a fixed physiological trait, it combines strong inter-writer distinctiveness with non-trivial intra-writer variability, making automated \emph{writer verification} -- deciding whether two handwriting samples share an author -- both valuable and challenging.

Verification differs from identification in that it is formulated as a one-to-one comparison between a probe sample and a claimed identity, rather than as a one-to-many search over a fixed set of enrolled identities. In writer-independent settings, the system must therefore generalize to writers not observed during training, making conventional closed-set identity classification unsuitable as the primary decision mechanism. This motivates pairwise and metric-learning formulations, with Siamese\cite{koch2015siamese} architectures representing a widely adopted paradigm for biometric verification, including face and signature verification.

Despite this maturity, applications to \emph{offline} (image-based) handwriting verification remain fragmented: prior studies typically fix a single backbone or a single training objective, rarely test cross-dataset generalization -- arguably the most operationally relevant scenario -- and seldom compare loss functions under controlled conditions, even though binary cross-entropy, contrastive, and triplet formulations encode substantially different notions of embedding-space structure.

We address these gaps with \textbf{HandVerify}, a modular Siamese framework enabling a systematic comparison along two axes: backbone encoder and training objective. We evaluate six CNN backbones across three families -- ResNet\cite{he2016resnet}-18/34, EfficientNet\cite{tan2019efficientnet}-B0/B1, MobileNetV3\cite{howard2019mobilenetv3}-Small/Large -- all ImageNet\cite{deng2009imagenet}-pretrained and adapted to grayscale input via channel-averaging of the first convolutional layer, preserving pretrained knowledge despite the RGB-to-grayscale modality shift. We compare three training paradigms: BCE, which frames verification as direct pairwise classification; Contrastive Loss\cite{hadsell2006contrastive}, which pulls genuine pairs together and separates impostor pairs by a margin; and Triplet Loss\cite{schroff2015facenet}, which enforces relative anchor-positive/anchor-negative separation. We further introduce two tailored training strategies: \emph{epoch-wise dynamic negative resampling}, addressing the combinatorial imbalance between scarce genuine pairs and abundant impostor pairs while exposing the model to diverse negatives; and \emph{selective backbone-layer freezing}, balancing retention of pretrained features against task-specific adaptation on limited data.

The framework is evaluated on IAM\cite{iam} and RIMES\cite{rimes} under four train$\to$test pairings -- two within-dataset (IAM$\to$IAM, RIMES$\to$RIMES) and two cross-dataset (IAM$\to$RIMES, RIMES$\to$IAM) -- for 72 trained (backbone, loss, pairing) configurations in total. This design lets us assess not only the best-performing combination, but also how consistently a loss function's advantage holds across backbones and domain shift. Performance is measured via EER, AUC, the decidability index $d'$, and fixed-operating-point classification metrics, complemented by paired Wilcoxon signed-rank testing for statistical consistency.

The paper is organized as follows. Section~\ref{sec:methodology} presents the architecture, loss formulations, and training strategies. Section~\ref{sec:experiments} describes datasets, preprocessing, and experimental protocol. Section~\ref{sec:results} reports the 72-configuration sweep, including same- and cross-dataset performance, significance analysis, and failure-case inspection. Section~\ref{sec:conclusion} summarizes findings, practical implications, limitations, and future directions.